\section{De Language Model a AI Agent: por qué una sola llamada no basta}

En el minuto 49, la clase pasa al tema de los agent. Esta transición no es un cambio de buzzword, sino un cambio en la unidad del sistema: un language model recibe tokens y genera tokens; un agent, además, tiene que mantener un state, leer la memory, elegir una action, observar el resultado en el entorno y seguir ejecutando según la retroalimentación. Una llamada al modelo puede ser un component del agent, pero no es un agent completo.

\subsection{Agent stack: qué más se necesita fuera del model}

La diapositiva de apertura enumera action, memory, communication y emergent architecture. La diapositiva Building AI Agents divide el problema en why y how: por qué se necesitan agent y cómo hacer que el modelo entre en interfaces reales.

\lecturefigure{slide-078.jpg}{Alcance del curso y líneas de investigación de la parte de AI Agent}{Diapositivas oficiales de CS25 V4, página 78}

\lecturefigure{slide-079.jpg}{Building AI Agents: marco del problema en términos de why y how}{Diapositivas oficiales de CS25 V4, página 79}

\begin{importantbox}{El ciclo mínimo de un agent}
Sea $s_t$ el estado del entorno; el modelo elige una action a partir de la observation $o_t$, la memory $m_t$ y el objetivo $g$:
\[
a_t=\pi_\theta(o_t,m_t,g),\qquad
s_{t+1}=\mathcal{E}(s_t,a_t),\qquad
m_{t+1}=U(m_t,o_t,a_t,s_{t+1}).
\]
$\pi_\theta$ es la policy, $\mathcal{E}$ es la transición del entorno y $U$ es la memory update. El agent tiene que procesar los resultados de sus acciones, no solo predecir el siguiente token.
\end{importantbox}

\lecturefigure{slide-080.jpg}{Una sola foundation-model call carece de estado a largo plazo y de retroalimentación de las acciones}{Diapositivas oficiales de CS25 V4, página 80}

\lecturefigure{slide-081.jpg}{La agent architecture agrega memory, tools y environment loop fuera del modelo}{Diapositivas oficiales de CS25 V4, página 81}

\begin{knowledgebox}{Términos clave: los seis componentes del agent stack}
\begin{tabularx}{\textwidth}{L{0.19\textwidth} X}
\toprule
Componente & Función \\
\midrule
Model/policy & Interpreta la observation y propone un plan o una action. \\
State & Guarda la tarea actual, las páginas, los archivos, las variables y el punto de ejecución. \\
Memory & Guarda información entre pasos o entre sesiones y permite recuperarla. \\
Tool/action layer & Asigna cada action estructurada a una API, un browser, un shell o un dispositivo. \\
Environment feedback & Devuelve éxito, errores, cambios en la página, resultados de pruebas u observation de sensores. \\
Verifier/guardrail & Verifica restricciones, permisos, resultados y condiciones de terminación. \\
\bottomrule
\end{tabularx}
\end{knowledgebox}

\teachervoice{El ponente explica de forma directa por qué se pasa a los agent: una sola foundation-model call no basta, porque las tareas reales también requieren memory, long context, personalization, actions, internet access y retroalimentación entre pasos.}

\subsection{Demos de Flight y driving-test: muestran action, no prueban confiabilidad}

Dos diapositivas de demo muestran a un agent temprano de MultiOn reservando un vuelo y completando el California online driving test. Muestran que el modelo puede convertir un objetivo en lenguaje natural en operaciones del navegador, pero un éxito único en un video público no permite inferir la tasa de éxito general, la seguridad de los pagos ni la robustez entre sitios web. Al leer la figura, primero identifica la observation, la action y la condición de término; después distingue la evidencia de demostración de la evidencia de despliegue.

\lecturefigure{slide-082.jpg}{Caso de reserva de vuelos con un AI agent mostrado en clase}{Diapositivas oficiales de CS25 V4, página 82}

\lecturefigure{slide-083.jpg}{Caso de agent para el online driving-test mostrado en clase}{Diapositivas oficiales de CS25 V4, página 83}

\begin{warningbox}{Cómo leer correctamente la demo evidence}
Una demo puede probar que «esta action trajectory ocurrió al menos una vez», pero no prueba la tasa de éxito promedio, la worst-case safety ni el despliegue sin supervisión. Para pasar de una capability demo a una engineering claim hacen falta un conjunto de tareas, el número de repeticiones, una clasificación de fallas, el alcance de los permisos, una estrategia de recuperación y una reproducción independiente.
\end{warningbox}

\teachervoice{El ponente llama a la demo del driving-test una de las primeras demostraciones de un real-world agent. Las notas conservan este hecho histórico, pero no lo presentan como la conclusión de que funcione de forma estable en cualquier sitio web, cuenta o dispositivo.}

\subsection{Por qué buscar una human-like interface}

El atractivo de un human-like agent es reutilizar las interfaces existentes: las páginas web, las aplicaciones de desktop y las mobile app ya están diseñadas para personas. Si el agent puede ver, hacer clic y escribir, no necesita esperar a que cada servicio ofrezca una API; además, puede aprender las preferencias del usuario a partir de sus demostraciones. Esa misma ventaja amplía el riesgo, porque el inicio de sesión, los pagos y los datos personales también quedan expuestos en el action space.

\lecturefigure{slide-084.jpg}{Un human-like agent reutiliza las interfaces, las preferencias y las demostraciones humanas}{Diapositivas oficiales de CS25 V4, página 84}

\begin{importantbox}{Interface coverage y safety boundary son dos caras de la misma moneda}
Un UI control más general aumenta la cobertura, pero debilita el schema, los permisos y los error code que proporciona una typed API. En el diseño conviene dar prioridad a las API con restricciones; solo cuando no exista una interfaz se amplía a UI control, y a las action de alto riesgo se les agregan confirmation, least privilege, audit log y rollback/recovery.
\end{importantbox}

\subsection{El autonomy level es una gradación de responsabilidad, no una categoría de marketing}

La clase toma prestada la clasificación L0--L5 de la conducción autónoma para hablar de agent: en los niveles bajos el control es humano y en los altos se retira gradualmente el human fallback. La analogía ayuda a explicar cómo cambia la responsabilidad, pero no sirve para clasificar directamente un automóvil o un producto de agent concreto. La siguiente figura debe leerse según «quién supervisa, quién puede tomar el control y quién responde cuando algo falla», no como una puntuación única de capacidad.

\lecturefigure{slide-085.jpg}{Analogía de cinco niveles, de human control a full autonomy}{Diapositivas oficiales de CS25 V4, página 85}

\begin{warningbox}{Los ejemplos de productos en clase son analogías informales}
El ponente asoció de palabra varios sistemas vehiculares con L4/L5. Las definiciones formales de SAE, el dominio de diseño operativo y el estado actual de los productos son mucho más rigurosos que la analogía de clase y cambian con el tiempo. La única conclusión didáctica que se conserva aquí es esta: retirar el human fallback eleva considerablemente los requisitos de verificación y de responsabilidad.
\end{warningbox}

\begin{knowledgebox}{Clasificación operativa de la agent autonomy}
\begin{tabularx}{\textwidth}{L{0.17\textwidth} X}
\toprule
Nivel & Una definición más adecuada para agent de software \\
\midrule
Assist & Ofrece sugerencias; no ejecuta action externas. \\
Approve-each-step & Cada action requiere la confirmación del usuario. \\
Bounded delegation & Ejecuta de forma automática dentro de un scope, un presupuesto y unas herramientas predefinidos. \\
Supervised autonomy & Puede ejecutar durante periodos largos, pero con monitor, interrupt y fallback en tiempo real. \\
Unsupervised autonomy & Completa tareas de alto impacto sin supervisión humana en línea; requiere los límites formales más estrictos y verificación independiente. \\
\bottomrule
\end{tabularx}
\end{knowledgebox}

\subsection{API y direct computer interaction}

Un agent tiene dos vías principales de action. La API ofrece parámetros estructurados, valores de retorno y permisos; la direct interaction opera las interfaces humanas mediante screen, keyboard, mouse o accessibility tree. La segunda tiene mayor cobertura, pero es más vulnerable al layout, la latency, los popup y la ambigüedad visual.

\lecturefigure{slide-087.jpg}{Dos vías: API control y direct computer interaction}{Diapositivas oficiales de CS25 V4, página 87}

\begin{importantbox}{Action schema}
Una action con capacidad de auditoría no debería ser solo texto libre, sino algo cercano a
\[
a_t=(\texttt{tool},\texttt{arguments},\texttt{precondition},\texttt{risk},\texttt{idempotency}).
\]
\texttt{precondition} describe el estado previo a la ejecución, \texttt{risk} determina si se requiere confirmación e \texttt{idempotency} indica si es seguro repetir la llamada. Las action estructuradas dan al verifier y a la recovery un objeto definido sobre el cual trabajar.
\end{importantbox}

\lecturefigure{slide-088.jpg}{Universal Action API y ejemplo de operación entre aplicaciones}{Diapositivas oficiales de CS25 V4, página 88}

\begin{warningbox}{«No se necesita API» no equivale a «la API no tiene valor»}
Un UI agent puede cubrir sistemas que no tienen API, pero una API estable suele ser más rápida, más barata y más fácil de autorizar y de probar. La arquitectura ideal es el capability routing: dar prioridad a las herramientas con restricciones y, cuando sea indispensable usar la UI, reducir el área visible, el action set y el credential scope.
\end{warningbox}

\subsection{Resumen de la sección}

El agent coloca al modelo de lenguaje dentro de un observation--action--feedback loop y le agrega state, memory, tool, verifier y recovery. Las demos de la clase mostraron la action capability, mientras que la autonomy y el UI control dejan al descubierto problemas de responsabilidad y de permisos. El verdadero límite del sistema no está en si «el modelo sabe hacer clic», sino en si las fallas pueden detectarse, acotarse y recuperarse.
